\section{Resumen ejecutivo}
% Una página con una panorámica del trabajo realizado que
% debería describir el objetivo, un párrafo describiendo la solución diseñada,
% aspectos clave y conclusiones. Este resumen debe contener todo lo necesario
% para familiarizarnos con el contenido de la memoria final.

Lo que se quiere conseguir a la hora de realizar este trabajo, es conseguir reunir en una única página web los lugares del planeta en los que se han conseguido certificar o validar avistamientos de ovnis. Mediante la utilización de una amplia gama de fuentes de datos heterogéneos  (CSV, APIs, www...) pertenecientes a diferentes satélites y empresas se quiere lograr trazar en un mapa global en tres dimensiones dichos datos de tal forma que se pueda ver la ubicación del avistamiento y los datos relevantes a ese suceso.
Estas amplias fuentes de datos cuentan con datos en diferentes formatos, por lo que se cuenta con datos auditivos, imágenes, video y demás. Además, se quiere utilizar los datos pertenecientes a distintas compañías, tanto públicas como privadas, con el fin de contrastar la información obtenida mediante diferentes fuentes sobre cada avistamiento.
Además, gracias al estado actual de la inteligencia artificial, se puede hacer uso de la misma con el fin de agilizar el trabajo, ya que por ejemplo, si se cuenta con audios muy largos, se puede utilizar esta tecnología para que seleccione las secciones más interesantes de ese audio.

Con el apoyo de una red interna Docker, Rust para el componente principal del \textit{backend} y Python como experto analista y extractor de datos, se busca alimentar una base de datos \textit{Mongo NoSQL} que ofrezca una prueba de concepto visual de la mano de React, de lo que puede llegar a ser una amplia red de avistamientos y fenómenos astronómicos que despierten el interés de aficionados y de la comunidad científica con el afán de seguir explorando el espacio.
